%-----------------------------------------------------------------------------------------------------------------------------------------------%
%	The MIT License (MIT)
%
%	Copyright (c) 2021 Jitin Nair
%
%	Permission is hereby granted, free of charge, to any person obtaining a copy
%	of this software and associated documentation files (the "Software"), to deal
%	in the Software without restriction, including without limitation the rights
%	to use, copy, modify, merge, publish, distribute, sublicense, and/or sell
%	copies of the Software, and to permit persons to whom the Software is
%	furnished to do so, subject to the following conditions:
%	
%	THE SOFTWARE IS PROVIDED "AS IS", WITHOUT WARRANTY OF ANY KIND, EXPRESS OR
%	IMPLIED, INCLUDING BUT NOT LIMITED TO THE WARRANTIES OF MERCHANTABILITY,
%	FITNESS FOR A PARTICULAR PURPOSE AND NONINFRINGEMENT. IN NO EVENT SHALL THE
%	AUTHORS OR COPYRIGHT HOLDERS BE LIABLE FOR ANY CLAIM, DAMAGES OR OTHER
%	LIABILITY, WHETHER IN AN ACTION OF CONTRACT, TORT OR OTHERWISE, ARISING FROM,
%	OUT OF OR IN CONNECTION WITH THE SOFTWARE OR THE USE OR OTHER DEALINGS IN
%	THE SOFTWARE.
%	
%
%-----------------------------------------------------------------------------------------------------------------------------------------------%

%----------------------------------------------------------------------------------------
%	DOCUMENT DEFINITION
%----------------------------------------------------------------------------------------

% article class because we want to fully customize the page and not use a cv template
\documentclass[a4paper,10pt]{article}

%----------------------------------------------------------------------------------------
%	FONT
%----------------------------------------------------------------------------------------

% % fontspec allows you to use TTF/OTF fonts directly
\usepackage{fontspec}
\defaultfontfeatures{Ligatures=TeX}
\setmainfont{Latin Modern Sans}


% % modified for ShareLaTeX use
% \setmainfont[
% SmallCapsFont = Fontin-SmallCaps.otf,
% BoldFont = Fontin-Bold.otf,
% ItalicFont = Fontin-Italic.otf
% ]
% {Fontin.otf}

%----------------------------------------------------------------------------------------
%	PACKAGES
%----------------------------------------------------------------------------------------
\usepackage{url}
\usepackage{parskip} 	

%other packages for formatting
\RequirePackage{color}
\RequirePackage{graphicx}
\usepackage[usenames,dvipsnames]{xcolor}
\usepackage[scale=0.9]{geometry}

%tabularx environment
\usepackage{tabularx}

%for lists within experience section
\usepackage{enumitem}

% centered version of 'X' col. type
\newcolumntype{C}{>{\centering\arraybackslash}X} 

%to prevent spillover of tabular into next pages
\usepackage{supertabular}
\usepackage{tabularx}
\newlength{\fullcollw}
\setlength{\fullcollw}{0.47\textwidth}

%custom \section
\usepackage{titlesec}				
\usepackage{multicol}
\usepackage{multirow}

%CV Sections inspired by: 
%http://stefano.italians.nl/archives/26
\titleformat{\section}{\Large\scshape\raggedright}{}{0em}{}[\titlerule]
\titlespacing{\section}{0pt}{10pt}{10pt}

%for publications
\usepackage[style=authoryear,sorting=ydnt, maxbibnames=5]{biblatex}


%Setup hyperref package, and colours for links
\usepackage[unicode, draft=false]{hyperref}
\definecolor{linkcolour}{rgb}{0,0.2,0.6}
\hypersetup{colorlinks,breaklinks,urlcolor=linkcolour,linkcolor=linkcolour}
\addbibresource{citations.bib}
\setlength\bibitemsep{1em}

%for social icons
\usepackage{fontawesome5}

%debug page outer frames
%\usepackage{showframe}

\DeclareNameAlias{sortname}{given-family}
\DeclareNameAlias{default}{given-family}


% job listing environments
\newenvironment{jobshort}[2]
    {
    \begin{tabularx}{\linewidth}{@{}l X r@{}}
    \textbf{#1} & \hfill &  #2 \\[3.75pt]
    \end{tabularx}
    }
    {
    }

\newenvironment{joblong}[2]
    {
    \begin{tabularx}{\linewidth}{@{}l X r@{}}
    \textbf{#1} & \hfill &  #2 \\[3.75pt]
    \end{tabularx}
    \begin{minipage}[t]{\linewidth}
    \begin{itemize}[nosep,after=\strut, leftmargin=1em, itemsep=3pt,label=--]
    }
    {
    \end{itemize}
    \end{minipage}    
    }

\newenvironment{job}[2]
{
	\begin{tabularx}{\linewidth}{@{}l X r@{}}
		\small
		\textbf{#1} & \hfill &  #2 \\[3.75pt]
	\end{tabularx}
	\footnotesize
	\setlength{\parskip}{3pt} % spacing between paragraphs
	\setlength{\parindent}{0pt} % no indentation
}
{
}


\usepackage{enumitem}

% Define compact list styles
\setlist[itemize,1]{nosep,after=\strut, leftmargin=2em, itemsep=3pt, label=•}
\setlist[itemize,2]{nosep,after=\strut, leftmargin=2em, itemsep=2pt, label=$\circ$}



%----------------------------------------------------------------------------------------
%	BEGIN DOCUMENT
%----------------------------------------------------------------------------------------
\begin{document}

% non-numbered pages
\pagestyle{empty} 

%----------------------------------------------------------------------------------------
%	TITLE
%----------------------------------------------------------------------------------------

% \begin{tabularx}{\linewidth}{ @{}X X@{} }
% \huge{Your Name}\vspace{2pt} & \hfill \emoji{incoming-envelope} email@email.com \\
% \raisebox{-0.05\height}\faGithub\ username \ | \
% \raisebox{-0.00\height}\faLinkedin\ username \ | \ \raisebox{-0.05\height}\faGlobe \ mysite.com  & \hfill \emoji{calling} number
% \end{tabularx}

\begin{tabularx}{\linewidth}{@{} C @{}}
\Huge{Dimitris Tsitsigkos} \\[7.5pt]
\Large{Software Engineer (PhD)} \\[7.5pt]

%\href{https://gr.linkedin.com/in/dimitris-tsitsigkos}{\raisebox{-0.05\height}\faLinkedin\ dtsitsigkos} \ $|$ \ 
%\href{https://dtsitsigkos.github.io/}{\raisebox{-0.05\height}\faGlobe \ dtsitsigkos.github.io} \ $|$ \  
%\href{https://dblp.org/pid/163/0537.html}{\raisebox{-0.05\height}\faDatabase\ DBLP}  \ $|$ \ 
%\href{https://scholar.google.com/citations?user=z_KVU0QAAAAJ}{\raisebox{-0.05\height}\faGraduationCap\ Google Scholar}\ $|$ \
%\href{https://github.com/dTsitsigkos}{\raisebox{-0.05\height}\faGithub\ dtsitsigkos} \ 

\href{mailto:tsitsigkosdim@gmail.com}{\raisebox{-0.05\height}\faEnvelope \ tsitsigkosdim@gmail.com} \ $|$ \
\href{tel:+30 6942951698}{\raisebox{-0.05\height}\faMobile \ +30 6942951698}
 \\
\end{tabularx}

%----------------------------------------------------------------------------------------
% EXPERIENCE SECTIONS
%----------------------------------------------------------------------------------------

%Interests/ Keywords/ Summary
%\section{Summary}
%\small Software Engineer (PhD) with experience building data management and parallel processing systems in research environments. Strong in performance-critical C++ development and in-memory indexing, with practical experience in distributed data processing and backend engineering using Java and Python. Focused on scalable architectures, efficient algorithms, and high-performance data systems.


\vspace{-0.7\baselineskip}
%Experience
\section{Work Experience}

\begin{job} {Postdoctoral researcher (Data Science and Engineering team)}{Feb 2025 – Present\phantom{Ja}}	
\textbf{Archimedes Research on AI, Data Science and Algorithms, Athens, Greece} 

\begin{itemize}

\item Designed and implemented high-performance, in-memory multi-dimensional and vector indices in C++, using parallelism and hardware-aware optimizations to improve query throughput and latency.

\item \underline{Technologies:} \textbf{C++, OpenMP, SIMD}
\end{itemize}

\end{job}


\begin{job}{PhD fellow }{Jan 2024 – Dec 2024}
	\textbf{Department of Computer Science \& Engineering, Univ. of Ioannina, Greece}
	\begin{itemize}
		
		%\item Implemented spatial index structures in C++ for spatial join, range queries, and k-NN, using parallel algorithms and memory-efficient data partitioning to scale on multi-core systems.
		
		%traditional and incremental 𝑘-NN, 𝜖-range, and two different
		%algorithms for 𝜖-distance joins
		
		\item Implemented spatial index structures in C++ for distance queries (k-NN, range, distance joins), improving query performance compared to existing approaches.
		
		
		\item Contributed to the design and development of a prototype distributed spatial data management system using MPI and OpenMP, reducing memory consumption by over 55\% while improving query performance.
		
		\item \underline{Technologies:} \textbf{C++, OpenMP, MPI}
		
	\end{itemize}
	
\end{job}

\vspace{-1\baselineskip}


\begin{job}{PhD fellow}{Jul 2019 – Dec 2024}
	\textbf{Information Management Systems Institute, Athena Research Center, Athens, Greece} 
	
	\begin{itemize}
	
		\item PhD Research: 
	
		\begin{itemize}
			
			\item Optimized B+-Tree Index Structures
				\begin{itemize}
			
				\item Designed and implemented BS-tree and CBS-tree, data-parallel B+-tree–based index structures optimized for in-memory query processing
			
				\item Achieved 2.5× higher query throughput vs B+-trees, learned indexes, and trie-based approaches
			
				\item Reduced memory consumption by up to 90\% using compressed layouts
			
				\item Designed for parallel execution and multi-core scalability
			
			\end{itemize}
		
			%\item Developed data-parallel, multithreaded BS‑tree and CBS‑tree (compressed variant) for in-memory indexing, achieving up to 2.5× faster single-threaded query throughput and up to 90\% memory savings compared to conventional B+‑trees, learned indexes, and trie-based structures, while demonstrating excellent data scalability, query scalability, and multi-core performance.
		\item Indexing and Partitioning for Non-Point Spatial Data
		
			\begin{itemize}
				\item Developed two-layer space-oriented partitioning for indexing non-point spatial objects
				
				\item Reduced spatial join cost by ~50\% compared to existing spatial indexing methods
				
				\item Achieved 2–3× faster range query performance vs R-trees, quad-trees and grid-based indexes
				
				\item Designed techniques supporting parallel spatial query processing
			\end{itemize}
			
			
		\item Parallel Spatial Join Processing
			\begin{itemize}
				\item Implemented a parallel spatial join algorithm based on the Partition-Based Spatial Merge (PBSM) algorithm
				
				\item Designed 1D and 2D partitioning strategies for efficient join processing
				
				\item Optimized sweep strategies to improve performance and scalability
				
				\item Enabled spatial joins over tens of millions of objects in under one second
				
			\end{itemize}
	
			%\item Implemented two-layer space-oriented partitioning for non-point spatial data, achieving ~50\% reduction in spatial join cost and 2–3× faster range query throughput compared to state-of-the-art spatial indexes (R‑trees, quad-trees, and grid-based methods), while supporting parallel processing and multi-core scalability.
			
			%\item Implemented a parallel in-memory spatial join framework by re-designing and tuning the PBSM (Partition-Based Spatial Merge) algorithm, including 1D/2D partitioning and sweep strategy optimization, enabling spatial joins of tens of millions of objects to complete in under one second while demonstrating excellent multi-core scalability.
	
		
		
	\end{itemize}
		
		\vspace{-1\baselineskip}
		
		\item Participated in European funding project:
		\begin{itemize}
		
			\item \href{https://www.more2020.eu/}{MORE}: Management of Real-time Energy data. 
				
			\begin{itemize}
				\item Built parallel continuous evaluation module in Java for sliding-window aggregations, enabling low-latency energy analytics.
				\item Implemented parallel versions of time-series analytics algorithms in Python using Dask to improve efficiency and scalability for large-scale data processing.
			
				%\item \underline{Technologies:} \textbf{Java, Python, Dask}.
			
				%\item Parallel sliding-window aggregation module in Java and distributed, parallel time-series analytics in Python (Dask).
			
			\end{itemize}
		\end{itemize}
		
		\vspace{-1\baselineskip}
		
		\item \underline{Technologies:} \textbf{C++, OpenMP, SIMD, Python, Dask}
		
	\end{itemize}
\end{job}

\vspace{-0.6\baselineskip}
\begin{job}{Software Engineer}{Apr 2017 – Nov 2018}
	\textbf{Hellenic Army Information Technology Support Center, Athens, Greece} 
	
	\begin{itemize}
		
		
		\item Maintained and extended Java backend applications for internal systems, improving performance and implementing new features.
		
		\item \underline{Technologies:} \textbf{Java, Oracle Database, JSF}.
	\end{itemize}
	
\end{job}


\begin{job}{Research Engineer}{Dec 2012 – Jun 2019}
	\textbf{Information Management Systems Institute, Athena Research Center, Athens, Greece} 
	
		
	Participated in several Greek and European funding projects, including:
	
	\begin{itemize}
		
		\item \href{https://amnesia.openaire.eu/}{Amnesia}: A platform for anonymizing relational, multi-dimensional, and hierarchical data. 
		
		%(Aug 2015 – Sep 2020)
		\begin{itemize}
			
			\item Implemented an open-source data anonymization platform for in-memory anonymization, supporting both web and desktop applications.
			
			\item Designed the architecture and full stack, with a Java backend providing RESTful APIs (Spring) and a web frontend using JavaScript and HTML.
			
			\item Integrated the platform into the  \href{https://www.openaire.eu/}{OpenAIRE} ecosystem, enabling global researchers, data providers, and companies to anonymize data efficiently.
			
			%\item \underline{Technologies:} \textbf{Java, Spring, RESTful Web Services, JavaScript}.
			
			%\item Built in-memory anonymization platform (Java/Spring/REST), widely used by researchers.
			
		\end{itemize}
		
		\vspace{-1\baselineskip}
		
		\item \href{http://www.modissense.gr/}{MoDisSENSE}: A Distributed Spatio-Temporal and Textual Processing Platform for Social. 
		
		%(Dec 2012 – Jul 2015)
		\begin{itemize}
			%\item Developed distributed data processing algorithms in Java using Hadoop and HBase to analyze large-scale location data, including point-of-interest discovery and trajectory reconstruction from GPS traces.
			%\item Built backend RESTful APIs and services with PostgreSQL for spatial data processing and recommendations.
			%\item \underline{Technologies:} \textbf{Java, RESTful Web Services, MapReduce, Hadoop, HBase, PostgreSQL}.
			
			
			\item Built distributed spatio-temporal processing (Java/Hadoop/HBase) for trajectory reconstruction and POI analytics with RESTful APIs and PostgreSQL.
			
		\end{itemize}
		
		\vspace{-1\baselineskip}
	
		\item \underline{Technologies:} \textbf{Java, RESTful Web Services, Spring, MapReduce, Hadoop, HBase, PostgreSQL}
	
	\end{itemize}
\end{job}

%\begin{job}{Software Engineer (Concurrent with MSc/PhD)}{Dec 2012 – Dec 2023}
%	\textbf{Information Management Systems Institute, Athens, Greece}
	
%	Participated in several \underline{Greek and European funding projects}, including:
%	\begin{itemize}
		
%		\item \href{https://www.more2020.eu/}{MORE}: Management of Real-time Energy data. 
		
		%(Oct 2020 – Dec 2023)		
%		\begin{itemize}
%			\item Built parallel continuous evaluation module in Java for sliding-window aggregations, enabling low-latency energy analytics.
%			\item Implemented parallel and distributed versions of time-series analytics algorithms in Python using Dask to improve efficiency and scalability for large-scale data processing.
			
			%\item \underline{Technologies:} \textbf{Java, Python, Dask}.
		
			%\item Parallel sliding-window aggregation module in Java and distributed, parallel time-series analytics in Python (Dask).
		
%		\end{itemize}
		
%		\vspace{-1\baselineskip}
		
%		\item \href{https://amnesia.openaire.eu/}{Amnesia}: A platform for anonymizing relational, multi-dimensional, and hierarchical data. 
		
		%(Aug 2015 – Sep 2020)
%		\begin{itemize}
			
%			\item Implemented an open-source data anonymization platform for in-memory anonymization, supporting both web and desktop applications.
			
%			\item Designed the architecture and full stack, with a Java backend providing RESTful APIs (Spring) and a web frontend using JavaScript and HTML.
			
%			\item Integrated the platform into the  \href{https://www.openaire.eu/}{OpenAIRE} ecosystem, making it accessible to a global community of researchers and data providers.
			
			%\item \underline{Technologies:} \textbf{Java, Spring, RESTful Web Services, JavaScript}.
			
			%\item Built in-memory anonymization platform (Java/Spring/REST), widely used by researchers.
			
%		\end{itemize}
		
%		\vspace{-1\baselineskip}
		
%		\item \href{http://www.modissense.gr/}{MoDisSENSE}: A Distributed Spatio-Temporal and Textual Processing Platform for Social. 
		
		%(Dec 2012 – Jul 2015)
%		\begin{itemize}
			%\item Developed distributed data processing algorithms in Java using Hadoop and HBase to analyze large-scale location data, including point-of-interest discovery and trajectory reconstruction from GPS traces.
			%\item Built backend RESTful APIs and services with PostgreSQL for spatial data processing and recommendations.
			%\item \underline{Technologies:} \textbf{Java, RESTful Web Services, MapReduce, Hadoop, HBase, PostgreSQL}.
			
			
%			\item Built distributed spatio-temporal processing (Java/Hadoop/HBase) for trajectory reconstruction and POI analytics with RESTful APIs and PostgreSQL.
			
%		\end{itemize}
		\vspace{-1\baselineskip}
		
%		\item \underline{Technologies:} \textbf{Java, Python, RESTful Web Services, Spring, MapReduce, Hadoop, Dask, HBase, PostgreSQL.}
		
%	\end{itemize}
%\end{job}

%\vspace{-1\baselineskip}



%----------------------------------------------------------------------------------------
%	EDUCATION
%----------------------------------------------------------------------------------------
\section{Education}


\begin{job}{PhD, Data Management}{Jul 2019 – Dec 2024}
	Department of Computer Science \& Engineering, University of Ioannina (UoI) in collaboration with Athena Research Center.
	
	\begin{itemize}
		\item Visiting Researcher (PhD Collaboration), Institute of Computer Science (Data Management research group), Johannes Gutenberg University Mainz, Germany (May 2022 – Jul 2022)
		
	\end{itemize}
	\vspace{-1\baselineskip}
	\underline{\textit{Thesis: In-memory Indexing for Parallel Processing of Single and Multi-Dimensional Queries.}}
	
	\underline{	Published papers in:} ICDE, VLDBJ, TKDE and SIG-SPATIAL
\end{job}

\vspace{0.2\baselineskip}


\begin{job}{MSc, Computing Systems: Software and Hardware, Computer Science}{Nov 2012 – Sep 2016}
	Department of Informatics and Telecommunications, National and Kapodistrian University of Athens (NKUA)
	
	\underline{\textit{Thesis: Complex Event Processing (CEP) for Intrusion Detection.}}	
\end{job}

\vspace{0.2\baselineskip}

\begin{job}{BSc, Computer Science}{Sep 2006 – Jun 2012}
	Department of Informatics and Telecommunications, National and Kapodistrian University of Athens (NKUA)
	
	\underline{\textit{Thesis: Clustering Wikipedia resources.}}
\end{job}





%----------------------------------------------------------------------------------------
%	Technical Skills
%----------------------------------------------------------------------------------------
\section{Technical Skills}
\begin{tabularx}{\linewidth}{@{}l X@{}}
	Languages: &  \normalsize{C, C++, Java, Python
	}\\
	Data Management:  &  \normalsize{MySQL, PostgreSQL, PostGIS, HBase}\\ 
	Parallel \& Distributed Systems: &  \normalsize{OpenMP, SIMD, MPI, Hadoop, Spark, Dask}\\
	Backend \& Web: &  \normalsize{Spring, RESTful Web Services, JavaScript, JSP, JSF
	}\\
	Operating Systems: &  \normalsize{Ubuntu, Microsoft Windows, macOS
	}\\
	Other: &  \normalsize{Data management, in-memory indexing \& optimization, parallel \& distributed Systems
	}\\   
\end{tabularx}


%----------------------------------------------------------------------------------------
%	Awards
%----------------------------------------------------------------------------------------

\section{Additional Experience \& Awards}

%\begin{job}{3rd place — %\href{https://databasecontest2024.athenarc.gr/}{Future of Database Programming Contest}, Athens} {Mar 2025}	

\begin{itemize}
	\item 3rd place — \href{https://databasecontest2024.athenarc.gr/}{Future of Database Programming Contest}, Athens (Mar 2025)
	\item External reviewer or member of program committee for about 8 conferences and journals on data management
	\item Published more than 10 peer-reviewed papers in international conferences and journals (h-index: 5)
	\item Volunteer: European Data Forum 2014, EDBT/ICDT 2023 Joint Conference, 6th ACM Europe Summer School on Data Science 2025, HDMS 2025
\end{itemize}
	
%\end{job}


%----------------------------------------------------------------------------------------
%	PUBLICATIONS
%----------------------------------------------------------------------------------------


%\section{Selected Publications}


%\begin{refsection}[citations.bib]
%	\setlength\bibitemsep{0.5em}
%	\nocite{*}
%	\printbibliography[heading=none]
%\end{refsection}


%----------------------------------------------------------------------------------------
%	Other
%----------------------------------------------------------------------------------------

%\vspace{-0.4 cm}
%\section{Other}
%\vspace{0.2\baselineskip}
%\begin{tabularx}{\linewidth}{@{}l X@{}}
%	Languages
%	&  \normalsize{Greek (native), English (Advanced)	}\\
	
%	Volunteer	&  \normalsize{European Data Forum 2014, EDBT/ICDT 2023 Joint Conference, 6th ACM Europe Summer School on Data Science 2025, HDMS 2025.
		
%	}\\ 
	
	
%\end{tabularx}

\section{Languages}
%\vspace{0.2\baselineskip}
\begin{itemize}
	\item Greek (native)
	\item English (Advanced)
\end{itemize}



	
\section{External Links}
\begin{itemize}
	\item LinkedIn: \texttt{https://gr.linkedin.com/in/dimitris\_tsitsigkos}
	\item Google Scholar: \texttt{https://scholar.google.com/citations?user=z\_KVU0QAAAAJ}
	\item DBLP: \texttt{https://dblp.org/pid/163/0537.html}
	\item GitHub: \texttt{https://github.com/dTsitsigkos}
	\item Personal website: \texttt{https://dtsitsigkos.github.io/}
\end{itemize}


\end{document}
